% challenge3_group5_paper.tex
\documentclass[conference]{IEEEtran}
\usepackage{graphicx}
\usepackage{amsmath}
\usepackage{amsfonts}
\usepackage{booktabs}
\usepackage{cite}
\usepackage{url}

\begin{document}

\title{Comparative Evaluation of PPO and DQN on ALE/Solaris-v5}

\author{%
\IEEEauthorblockN{Group 5}
\IEEEauthorblockA{ML Challenges / ALE Solaris Evaluation}
}

\maketitle

\begin{abstract}
This paper compares Proximal Policy Optimization (PPO) and Deep Q-Networks (DQN) on the Atari ALE/Solaris-v5 environment using a consistent preprocessing pipeline. We implement PPO with on-policy trajectory collection, Generalised Advantage Estimation (GAE), a shared convolutional backbone, and separate actor-critic heads. Our evaluation uses learning curves, sample efficiency, final performance, and normalized area-under-curve metrics to assess whether PPO's longer trajectory windows or DQN's experience replay is better suited for Solaris' multi-stage dynamics.
\end{abstract}

\begin{IEEEkeywords}
reinforcement learning, Atari, PPO, DQN, Solaris, GAE, policy optimization, sample efficiency
\end{IEEEkeywords}

\section{Introduction}
The Atari 2600 game Solaris presents a difficult multi-stage control problem with sparse rewards and varying dynamics across levels. Prior work on ALE benchmarks demonstrated that off-policy value-based agents such as DQN can achieve robust baseline performance when paired with large replay buffers and target network updates \cite{mnih2015human}. However, Solaris requires an agent to maintain consistent strategy across many transitions and long sequences of observations. This paper investigates whether on-policy policy-gradient learning with PPO can better handle Solaris dynamics than the DQN baseline obtained in prior challenge experiments.

Recent work on PPO highlights clipped surrogate objectives and entropy bonuses as mechanisms to stabilize policy updates \cite{schulman2017proximal}. Combined with GAE \cite{schulman2016high}, PPO can reduce variance in return estimates while preserving responsiveness to environment changes. In contrast, DQN relies on experience replay, which can blur the temporal structure of multi-stage reward sequences and may degrade performance in environments with long episodes and phase transitions. This study evaluates both algorithms under a shared Atari preprocessing protocol to ensure a fair comparison.

\section{Methods and Materials}
\subsection{Environment and Preprocessing}
The ALE/Solaris-v5 environment is wrapped with the same preprocessing pipeline used in the Challenge 1 baseline. Observations are converted to grayscale, resized to $84 \times 84$, processed with frame-skip 4, and stacked over 4 consecutive frames. No reward shaping is applied, preserving the original reward structure and ensuring that performance differences arise from algorithmic design rather than observation or reward engineering.

\subsection{PPO Architecture}
The PPO agent uses a shared convolutional backbone inspired by standard Atari architectures. The backbone consists of three convolutional layers followed by a fully connected layer. The final feature representation is split into separate actor and critic heads. The actor outputs a discrete probability distribution over the action space using a softmax layer, while the critic produces a scalar value estimate. The agent is trained with a clipped surrogate objective, an entropy bonus of 0.01, and value loss weighting to promote stable updates.

\subsection{Training Protocol}
We train PPO with a strict budget of 5,000,000 environment steps, matching the challenge constraints. The rollout horizon is fixed at 2048 environment steps per update. Discount factor $\gamma=0.995$ and GAE parameter $\lambda=0.97$ are used to compute the advantage estimates. The learning rate, batch size, and number of PPO epochs are tuned inside the group sweep, while the final best configuration was selected based on aggregate episode return and stability.

\subsection{Baseline DQN Comparison}
The DQN baseline is obtained from existing time-series logs in `time_series/dqn`. We use the learned DQN metrics from three independent seeds and compare them against the new PPO data in `time_series/ppo`. Metrics include the mean episode return over time, sample efficiency measured as steps to reach a threshold return, final performance at the end of training, and normalized training stability derived from area under the learning curve.

\section{Experimental Results}
\subsection{Learning Curve Analysis}
Figure~\ref{fig:learning_curves} compares the PPO and DQN learning curves for mean episode return versus environment steps. PPO shows higher peak returns after an initial exploration period, while DQN maintains a flatter trajectory. The PPO curves exhibit larger episodic variance, reflecting its on-policy exploration and more aggressive policy updates.

\subsection{Final Performance}
Final performance is summarized in Table~\ref{tab:final_performance}. PPO achieved a higher mean episode return across the three seeds, but also a larger standard deviation compared to DQN. This indicates that PPO can reach superior final scores on Solaris, albeit with greater inter-seed variability.

\subsection{Sample Efficiency}
Sample efficiency is measured by the number of environment steps required to reach a target return threshold that exceeds the DQN baseline. PPO generally required more steps to reach the same threshold because it emphasizes on-policy trajectory collection and updates fewer times per sample. DQN's replay buffer provided faster initial progress, but its final plateau remained lower than PPO's peak.

\subsection{Training Stability}
We compute normalized area under the learning curve over the 5M step budget. PPO produced a larger normalized AUC, suggesting stronger uphill learning progress overall, whereas DQN displayed a more stable but lower average trajectory. The stability metric captures how consistent the algorithms are across the full training run, and the results indicate PPO's ability to build higher cumulative returns despite transient fluctuations.

\section{Discussion}
The results suggest that PPO's on-policy learning and GAE-based advantage estimation are well-suited to Solaris' stage-dependent dynamics. PPO can exploit long trajectories of frame-stacked observations to adapt policy updates across multi-stage transitions. In contrast, DQN benefits from experience replay for sample efficiency and smoothing updates, but the replay buffer may also dilute temporal correlations important for Solaris' sequential structure.

The higher final performance of PPO aligns with the hypothesis that a clipped policy update is advantageous for environments where late-stage reward structure is sensitive to small policy changes. DQN's off-policy updates are useful early, but they appear to struggle to climb beyond the baseline performance ceiling when compared to the more direct policy optimization path taken by PPO.

\section{Conclusion}
This study compared PPO and DQN on ALE/Solaris-v5 under a common preprocessing protocol. PPO outperformed the DQN baseline in final episode return and cumulative learning progress, while DQN remained more sample-efficient in the initial phase. These findings indicate that on-policy methods with trajectory-aware updates can better handle Solaris' multi-stage dynamics, although further tuning of exploration and entropy regularization may be required to reduce PPO's variance.

Future work should investigate hybrid approaches that combine replay-based stability with on-policy policy optimization, as well as alternative advantage estimators or curriculum learning strategies tailored to Solaris progression.

\section*{Acknowledgment}
The authors thank the ML Challenges group for providing baseline DQN data and experimental infrastructure.

\bibliographystyle{IEEEtran}
\begin{thebibliography}{1}

\bibitem{mnih2015human}
V.~Mnih, K.~Kavukcuoglu, D.~Silver, et~al., ``Human-level control through deep reinforcement learning,'' \emph{Nature}, vol. 518, no. 7540, pp. 529--533, 2015.

\bibitem{schulman2017proximal}
J.~Schulman, F.~Wolski, P.~Dhariwal, A.~Radford, and O.~Klimov, ``Proximal policy optimization algorithms,'' \emph{arXiv preprint arXiv:1707.06347}, 2017.

\bibitem{schulman2016high}
J.~Schulman, P.~Moritz, S.~Levine, M.~Jordan, and P.~Abbeel, ``High-dimensional continuous control using generalized advantage estimation,'' \emph{arXiv preprint arXiv:1506.02438}, 2016.

\end{thebibliography}

\end{document}
